\section{The analytic term}
\label{s-relaxed-excess}

Define the Chebyshev-normalized Mertens error
\begin{equation}
\label{e-mertens-error}
    E_\theta(x)
    =
    -\sum_{p\leq x}\log\left(1-\frac1p\right)
    -\gamma-\log\log\theta(x).
\end{equation}
The relaxed problem differs from this classical term by less than the
integrality gap.

\paragraph{Relaxed value.}
We first show that
\begin{equation}
\label{e-relaxed-value}
    g_{\reals}(m)
    =
    E_\theta(p_m)
    +
    O\left(
        \frac1{\sqrt{p_m}\log^{3/2}p_m}
    \right).
\end{equation}
Put \(y=\sqrt{1+T^\star}\). The response
\eqref{e-relaxed-response} gives
\[
    p_i-p_i^{-r_i^\star}
    =
    \frac{p_iT^\star}{1+T^\star}
\]
for \(p_i\leq y\), while \(r_i^\star=1\) for \(p_i>y\).
Substitution in \eqref{e-objective} gives
\[
\begin{aligned}
    g_{\reals}(m)
    &=
    \pi(y)\log\frac{T^\star}{1+T^\star}
    -\sum_{p\leq p_m}\log\left(1-\frac1p\right)\\
    &\quad
    +\sum_{y<p\leq p_m}\log\left(1-\frac1{p^2}\right)
    -\gamma-\log\log S^\star.
\end{aligned}
\]
The prime number theorem and partial summation give
\[
    \pi(y)\left|\log\frac{T^\star}{1+T^\star}\right|
    +
    \sum_{p>y}\left|\log\left(1-\frac1{p^2}\right)\right|
    =
    O\left(\frac1{\sqrt{T^\star}\log T^\star}\right).
\]
Proposition~\ref{prop-relaxed-scale} and
\eqref{e-relaxed-displacement} also give
\[
    \log\log S^\star-\log\log\theta(p_m)
    =
    O\left(\frac1{\sqrt{T^\star}\log T^\star}\right).
\]
The result follows from \(T^\star\sim p_m\log p_m\).

We now use Nicolas's results to evaluate \(E_\theta(p_m)\).

\subsection{Under the Riemann hypothesis}
\label{s-rh-branch}

Assume the Riemann hypothesis, and define
\[
    \mathcal H(t)
    =
    \sum_\rho
    \frac{e^{(\rho-1/2)t}}{\rho(\rho-1)},
    \qquad
    \tau=2+\gamma-\log(4\pi),
\]
where the sum runs over the nontrivial zeros of the zeta function, with
multiplicity. The series is absolutely convergent because its terms are
\(O(|\operatorname{Im}\rho|^{-2})\).

Nicolas's integrated explicit formula
\cite[Theorem~3(b), final formula in its proof]
{nicolas-1983-petites-valeurs}, translated to the normalization
\eqref{e-mertens-error}, is
\begin{equation}
\label{e-nicolas-integrated}
    E_\theta(x)
    =
    \frac{2-\mathcal H(\log x)}
         {\sqrt x\log x}
    +
    O\left(\frac{1}{\sqrt x\log^2x}\right).
\end{equation}
The same proof records the classical identity
\[
    \sum_\rho\frac{1}{\rho(1-\rho)}
    =
    2+\gamma-\log(4\pi).
\]
On the critical line, \(\rho(1-\rho)=|\rho|^2\). Therefore
\begin{equation}
\label{e-H-bound}
    |\mathcal H(t)|\leq\tau.
\end{equation}

\begin{theorem}[Relaxed excess under the Riemann hypothesis]
\label{thm-rh-relaxed}
Assume the Riemann hypothesis. Then
\[
    \sqrt{p_m}\log p_m\,g_{\reals}(m)
    =
    2-\mathcal H(\log p_m)
    +O\left(\frac1{\sqrt{\log p_m}}\right).
\]
In particular,
\[
    2-\tau+o(1)
    \leq
    \sqrt{p_m}\log p_m\,g_{\reals}(m)
    \leq
    2+\tau+o(1),
\]
and \(g_{\reals}(m)>0\) for all sufficiently large \(m\).
\end{theorem}

\begin{proof}
Apply \eqref{e-nicolas-integrated} at \(x=p_m\), and use
\eqref{e-relaxed-value}. The error in \eqref{e-relaxed-value}, after
multiplication by \(\sqrt{p_m}\log p_m\), is
\(O(1/\sqrt{\log p_m})\). The bounds follow from
\eqref{e-H-bound}. The lower one is eventually positive because
\(2-\tau>0\).
\end{proof}

The bound describes a band containing the normalized relaxed excess. It
does not assert that the full band is its limit set or that either
endpoint is attained along the primes.

\subsection{Off the Riemann hypothesis}
\label{s-off-rh-branch}

We use \(f=\Omega_\pm(g)\) when \(f/g\) has a positive limsup and a
negative liminf.

\begin{theorem}[Off-RH excursions]
\label{thm-off-rh-relaxed}
Suppose the Riemann hypothesis is false, and let
\[
    \Theta
    =
    \sup\{\operatorname{Re}\rho\mid\zeta(\rho)=0\}.
\]
For every \(0<\eps<\Theta-1/2\),
\[
    g_{\reals}(m)
    =
    \Omega_\pm\left(p_m^{-1/2+\eps}\right).
\]
Consequently,
\[
\begin{aligned}
    \limsup_{m\to\infty}
    \sqrt{p_m}\log p_m\,g_{\reals}(m)
    &=+\infty,\\
    \liminf_{m\to\infty}
    \sqrt{p_m}\log p_m\,g_{\reals}(m)
    &=-\infty.
\end{aligned}
\]
\end{theorem}

\begin{proof}
Nicolas
\cite[Theorem~3(c)]{nicolas-1983-petites-valeurs} proves, in the
normalization \eqref{e-mertens-error}, that
\[
    E_\theta(x)=\Omega_\pm(x^{-b})
    \qquad
    (1-\Theta<b<1/2).
\]
The function \(E_\theta\) is constant between consecutive primes, and
the prime number theorem makes the preceding prime asymptotic to \(x\).
The two real sequences in Nicolas's result therefore give two prime
subsequences with the same orders and opposite signs.

By \eqref{e-relaxed-value},
\[
    g_{\reals}(m)
    =
    E_\theta(p_m)
    +
    O\left(
        \frac{1}{\sqrt{p_m}\log^{3/2}p_m}
    \right).
\]
The error is \(o(p_m^{-b})\). Taking \(b=1/2-\eps\) proves the first
claim, and the normalized divergence follows.
\end{proof}
